\documentclass[12pt]{article}

\usepackage{amsmath, amssymb}

\usepackage[margin=1in]{geometry}

\begin{document}

% =========================
% Problem 1
% =========================
\section*{Problem 1}
\begin{itemize}
    \item[(a)]
        For all $n \in \mathbb{N}$, let $P(n)$ be the statement that $x_{n+1} > x_n$.\\

\textbf{Base case.}  
When $n = 1$, we have
\[
x_1 = 0, \qquad
x_2 = \sqrt{x_1 + 10} = \sqrt{0 + 10} = \sqrt{10}.
\]
Since $\sqrt{10} > 0 = x_1$, the statement $P(1)$ is true.

\medskip

\textbf{Inductive hypothesis.}  
Assume that $P(n)$ is true for some $n \in \mathbb{N}$, that is,
\[
x_{n+1} > x_n.
\]

\medskip

\textbf{Inductive step.}  
Adding $10$ to both sides gives
\[
x_{n+1} + 10 > x_n + 10.
\]
Taking square roots of both sides,
\[
\sqrt{x_{n+1} + 10} > \sqrt{x_n + 10}.
\]
By the definition of the sequence,
\[
x_{n+2} > x_{n+1},
\]
so $P(n+1)$ is true.

\medskip

\textbf{Conclusion.}  
Since $P(1)$ is true and $P(n) \Rightarrow P(n+1)$, it follows that
\[
P(n) \text{ is true for all } n \in \mathbb{N}.
\]
    \item[(b)]


Let $P(n)$ be the statement that $x_{n+1} < x_n$ for all $n \in \mathbb{N}$.

\medskip

\textbf{Base case.}  
When $n = 1$, we have
\[
x_1 = 10 = 2 \cdot 5,
\qquad
x_2 = \sqrt{x_1 + 10} = \sqrt{20} = 2\sqrt{5}.
\]
Since $2\sqrt{5} < 10$, we have $x_2 < x_1$, so $P(1)$ is true.

\medskip

\textbf{Inductive hypothesis.}  
Assume that $P(n)$ is true for some $n \in \mathbb{N}$, that is,
\[
x_{n+1} < x_n.
\]

\medskip

\textbf{Inductive step.}  
Adding $10$ to both sides gives
\[
x_{n+1} + 10 < x_n + 10.
\]
Taking square roots of both sides,
\[
\sqrt{x_{n+1} + 10} < \sqrt{x_n + 10}.
\]
By the definition of the sequence,
\[
x_{n+2} < x_{n+1},
\]
so $P(n+1)$ is true.

\medskip

\textbf{Conclusion.}  
Since $P(1)$ is true and $P(n) \Rightarrow P(n+1)$, it follows that
\[
P(n) \text{ is true for all } n \in \mathbb{N}.
\]
\end{itemize}

% =========================
% Problem 2
% =========================
\section*{Problem 2}

Consider the polynomial
\[
f(x) = 10x^4 + 7x^3 - 4x^2 + 14x + 3.
\]

By the Rational Root Theorem, any rational root has the form
\[
x = \frac{a}{b},
\]
where $\gcd(a,b)=1$, $a \mid 3$, and $b \mid 10$.

\medskip

\textbf{Divisors.}
\[
\begin{aligned}
\text{Divisors of } 3 &: \quad \pm 1, \pm 3, \\
\text{Divisors of } 10 &: \quad \pm 1, \pm 2, \pm 5, \pm 10.
\end{aligned}
\]

\medskip

\textbf{Possible rational roots.}

Forming all reduced fractions $\frac{a}{b}$ gives
\[
\pm 1, \pm 3, \pm \frac{1}{2}, \pm \frac{3}{2},
\pm \frac{1}{5}, \pm \frac{3}{5},
\pm \frac{1}{10}, \pm \frac{3}{10}.
\]

\medskip

\textbf{Evaluation of all candidates.}

\[
\begin{array}{c|c}
x & f(x) \\ \hline
1 & 30 \\
-1 & -12 \\
3 & 900 \\
-3 & 576 \\
\frac{1}{2} & \frac{225}{8} \\
-\frac{1}{2} & -\frac{27}{8} \\
\frac{3}{2} & 192 \\
-\frac{3}{2} & 0 \\
\frac{1}{5} & \frac{336}{125} \\
-\frac{1}{5} & 0 \\
\frac{3}{5} & \frac{144}{25} \\
-\frac{3}{5} & -\frac{96}{25} \\
\frac{1}{10} & \frac{147}{1000} \\
-\frac{1}{10} & -\frac{63}{1000} \\
\frac{3}{10} & \frac{459}{1000} \\
-\frac{3}{10} & -\frac{231}{1000}
\end{array}
\]

\medskip

Since
\[
f\!\left(-\frac{1}{5}\right) = 0
\quad \text{and} \quad
f\!\left(-\frac{3}{2}\right) = 0,
\]
and all other candidates give nonzero values, the rational roots are
\[
\boxed{x = -\frac{1}{5} \quad \text{and} \quad x = -\frac{3}{2}}.
\]

% =========================
% Problem 3
% =========================
\section*{Problem 3}
\begin{itemize}
    \item[a)]

Consider the expression
\[
\sqrt{14 + 6\sqrt{5}} - \sqrt{6 + 2\sqrt{5}}.
\]

We rewrite each radical as a perfect square.

\[
14 + 6\sqrt{5}
= 9 + 5 + 2 \cdot 3 \sqrt{5}
= (3 + \sqrt{5})^2.
\]

\[
6 + 2\sqrt{5}
= 1 + 5 + 2 \cdot 1 \sqrt{5}
= (1 + \sqrt{5})^2.
\]

Therefore,
\[
\sqrt{14 + 6\sqrt{5}} - \sqrt{6 + 2\sqrt{5}}
= \sqrt{(3 + \sqrt{5})^2} - \sqrt{(1 + \sqrt{5})^2}.
\]

Since both quantities are positive, this simplifies to
\[
(3 + \sqrt{5}) - (1 + \sqrt{5}) = 2.
\]

Thus, the value of the expression is
\[
\boxed{2},
\]
which is rational.

    \item[b)]

Suppose that
\[
r = \frac{\log 2}{\log 3}
\]
is rational. Then there exist integers $p, q > 0$ such that
\[
r = \frac{p}{q}.
\]

Thus,
\[
\frac{\log 2}{\log 3} = \frac{p}{q}.
\]
Multiplying both sides by $\log 3$, we obtain
\[
\log 2 = \frac{p}{q} \log 3.
\]
Multiplying both sides by $q$ gives
\[
q \log 2 = p \log 3.
\]

Exponentiating both sides with base $10$, we have
\[
10^{q \log 2} = 10^{p \log 3}.
\]
Using the identity $10^{\log a} = a$, this simplifies to
\[
2^q = 3^p.
\]

However, $2^q$ is a power of an even number, while $3^p$ is a power of an odd number.  
Since no power of an even integer can equal a power of an odd integer, this is a contradiction.

Therefore, our assumption was false, and
\[
\boxed{\frac{\log 2}{\log 3} \text{ is irrational}.}
\]
\end{itemize}

% =========================
% Problem 4
% =========================
\section*{Problem 4}

\[
|x-2| + |x-4| < x + 1.
\]

The critical points occur when the expressions inside the absolute values are zero:
\[
x-2 = 0 \quad \text{and} \quad x-4 = 0,
\]
so the critical points are $x = 2$ and $x = 4$.\\

We consider three cases.

\medskip

\textbf{Case 1: $x < 2$.}

In this case, $|x-2| = 2-x$ and $|x-4| = 4-x$.  
Thus,
\[
(2-x) + (4-x) < x + 1.
\]
\[
6 - 2x < x + 1,
\]
\[
-3x < -5,
\]
\[
x > \frac{5}{3}.
\]
Hence,
\[
x \in \left(\frac{5}{3},\,2\right).
\]

\medskip

\textbf{Case 2: $2 \le x < 4$.}

Here, $|x-2| = x-2$ and $|x-4| = 4-x$.  
Thus,
\[
(x-2) + (4-x) < x + 1,
\]
\[
2 < x + 1,
\]
\[
x > 1.
\]
Since this holds for all $x$ in the interval, we have
\[
x \in [2,\,4).
\]

\medskip

\textbf{Case 3: $x \ge 4$.}

In this case, $|x-2| = x-2$ and $|x-4| = x-4$.  
Thus,
\[
(x-2) + (x-4) < x + 1,
\]
\[
2x - 6 < x + 1,
\]
\[
x < 7.
\]
Hence,
\[
x \in [4,\,7).
\]

\medskip

\textbf{Conclusion.}

Combining all three cases, the solution set is
\[
\boxed{\left(\frac{5}{3},\,7\right)}.
\]

% =========================
% Problem 5
% =========================
\section*{Problem 5}

\begin{itemize}

\item[(a)] 
To show that the set
\[
S = \left\{ \frac{m}{n} : m,n \in \mathbb{N},\ m < 3n \right\}
\]
is bounded, we show that it is bounded above and bounded below.

\medskip

\textbf{Bounded below.}  
Since $m,n \in \mathbb{N}$, we have $m \ge 1$ and $n \ge 1$, so
\[
\frac{m}{n} > 0
\]
for every element of $S$. Hence, $0$ is a lower bound for $S$, and $S$ is bounded below.

\medskip

\textbf{Bounded above.}  
Since $m < 3n$, dividing both sides by $n > 0$ gives
\[
\frac{m}{n} < 3.
\]
Thus, every element of $S$ is less than $3$, so $3$ is an upper bound for $S$, and $S$ is bounded above.

\medskip

\textbf{Conclusion.}  
Since $S$ is bounded both above and below, $S$ is bounded.

\medskip

\item[(b)]
\textbf{Finding $\inf(S)$ and $\sup(S)$.}

\medskip

\textbf{Infimum.}  
From part (a), since $m,n \in \mathbb{N}$, we have $m \ge 1$ and $n \ge 1$, so
\[
\frac{m}{n} > 0
\]
for all $\frac{m}{n} \in S$. Hence, $0$ is a lower bound of $S$.

\medskip

To show that $0$ is the greatest lower bound, let $\varepsilon > 0$ be arbitrary.
Choose $m = 1$ and let $n \in \mathbb{N}$ be sufficiently large. Then
\[
\frac{m}{n} = \frac{1}{n} \in S,
\]
since $1 < 3n$ for all $n \in \mathbb{N}$.

By the Archimedean Axiom Corollary, there exists $n \in \mathbb{N}$ such that
\[
\frac{1}{n} < \varepsilon.
\]
Thus, there exists $s \in S$ such that
\[
0 < s < \varepsilon.
\]

\medskip

Since for every $\varepsilon > 0$ there exists $s \in S$ with
\[
s < 0 + \varepsilon,
\]
no number greater than $0$ can be a lower bound of $S$. Therefore,
\[
\inf(S) = 0.
\]

\medskip

\textbf{Supremum.}  
From part (a), since $m < 3n$, dividing both sides by $n > 0$ gives
\[
\frac{m}{n} < 3
\]
for all $\frac{m}{n} \in S$. Hence, $3$ is an upper bound of $S$.

\medskip

To show that $3$ is the least upper bound, let $\varepsilon > 0$ be arbitrary.
Choose $m = 3n - 1$ with $n \in \mathbb{N}$. Then $m < 3n$, and
\[
s = \frac{m}{n} = \frac{3n - 1}{n} = 3 - \frac{1}{n} \in S.
\]

By the Archimedean Axiom Corollary, there exists $n \in \mathbb{N}$ such that
\[
\frac{1}{n} < \varepsilon.
\]
Thus,
\[
s = 3 - \frac{1}{n} > 3 - \varepsilon.
\]

\medskip

Since for every $\varepsilon > 0$ there exists $s \in S$ such that
\[
3 - \varepsilon < s < 3,
\]
no number smaller than $3$ can be an upper bound of $S$. Therefore,
\[
\sup(S) = 3.
\]

\end{itemize}

\end{document}
